\documentclass[11pt,a4paper]{article}
\usepackage[italian]{babel}
\usepackage[margin=2.5cm]{geometry}
\usepackage{parskip}
\usepackage{amsmath, amssymb, amsthm}
\usepackage{booktabs}
\usepackage{array}
\usepackage{xcolor}
\usepackage{needspace}
\usepackage{enumitem}
\usepackage{listings}
\usepackage{tikz}
\usepackage[most]{tcolorbox}
\usepackage[hidelinks]{hyperref}
\usetikzlibrary{decorations.pathreplacing, shapes.geometric, arrows.meta, positioning}

% Niente vedove/orfane: i paragrafi non lasciano righe isolate a fine/inizio pagina
\widowpenalty=10000
\clubpenalty=10000

% Elenchi più compatti
\setlist{itemsep=2pt, parsep=0pt, topsep=4pt, partopsep=0pt}

% --- Riquadri con bordo nero, senza numero (si spezzano tra due pagine) ---
% Uso: \begin{definizione} ... \end{definizione}
%      \begin{teorema}[Nome] ... \end{teorema}  (il nome è facoltativo)
\tcbset{nota/.style={enhanced jigsaw, breakable, colback=white,
  colframe=black, boxrule=0.5pt, sharp corners}}
\NewTColorBox{definizione}{}{nota,
  before upper={\textbf{Definizione.}\ }}
\NewTColorBox{teorema}{o}{nota,
  before upper={\textbf{Teorema\IfValueT{#1}{ (#1)}.}\ }}
\NewTColorBox{proposizione}{o}{nota,
  before upper={\textbf{Proposizione\IfValueT{#1}{ (#1)}.}\ }}
\NewTColorBox{proprieta}{o}{nota,
  before upper={\textbf{Proprietà\IfValueT{#1}{ (#1)}.}\ }}
\NewTColorBox{assioma}{o}{nota, boxrule=1pt,
  before upper={\textbf{Assioma\IfValueT{#1}{ (#1)}.}\ }}

% --- Ambienti semplici (senza riquadro) ---
% Il titolo sta in cima al blocco, anche se il contenuto inizia con un riquadro o
% con codice; e non resta mai da solo in fondo a una pagina.
% Uso: \begin{esempio}[Titolo facoltativo] ... \end{esempio}
\newcommand{\newrunin}[2]{%
  \NewDocumentEnvironment{#1}{o}{%
    \par\Needspace{4\baselineskip}\addvspace{\medskipamount}\noindent
    \textbf{#2}\IfValueT{##1}{ (##1)}\textbf{.}\ \ignorespaces}%
  {\par\addvspace{\medskipamount}}}
\newrunin{esempio}{Esempio}
\newrunin{esercizio}{Esercizio}
\newrunin{osservazione}{Osservazione}

% V e F nelle tavole di verità
\newcommand{\V}{\textsf{V}}
\newcommand{\F}{\textsf{F}}

% --- Codice C ---
% Uso: \begin{codice} ... \end{codice}      codice C numerato
%      \begin{comando} ... \end{comando}    comandi da terminale
%      \begin{risultato} ... \end{risultato} output di un programma
%      \lstinline|printf("ciao");|           codice dentro al testo
\lstdefinestyle{c}{
  language=C,
  basicstyle=\ttfamily\small,
  keywordstyle=\color{blue}\bfseries,
  commentstyle=\color{gray}\itshape,
  stringstyle=\color{red!70!black},
  numbers=left,
  numberstyle=\tiny\color{gray},
  numbersep=8pt,
  columns=flexible,
  keepspaces=true,
  breaklines=true,
  showstringspaces=false,
  tabsize=4,
  morekeywords={include, define},
  % lettere accentate nei commenti
  literate={à}{{\`a}}1 {è}{{\`e}}1 {é}{{\'e}}1 {ì}{{\`i}}1
           {ò}{{\`o}}1 {ù}{{\`u}}1 {È}{{\`E}}1 {À}{{\`A}}1
}
\lstset{style=c}
\newtcblisting{codice}{listing only, nota, left=6mm,
  listing options={style=c}}
\newtcblisting{comando}{listing only, nota,
  listing options={basicstyle=\ttfamily\small, numbers=none, language={},
    columns=flexible, keepspaces=true, breaklines=true}}
\newtcblisting{risultato}{listing only, enhanced jigsaw, breakable,
  colback=black!5, colframe=black, boxrule=0.5pt, sharp corners,
  listing options={basicstyle=\ttfamily\small, numbers=none, language={},
    columns=flexible, keepspaces=true, breaklines=true}}

% --- Diagrammi di flusso (simboli standard) ---
\tikzset{
  terminale/.style = {ellipse, draw, minimum width=2.5cm, minimum height=0.9cm},
  io/.style        = {trapezium, trapezium left angle=70, trapezium right angle=110,
                      draw, minimum width=2.5cm, minimum height=0.9cm},
  processo/.style  = {rectangle, draw, minimum width=2.5cm, minimum height=0.9cm},
  decisione/.style = {diamond, draw, aspect=2, minimum width=2.5cm},
  freccia/.style   = {thick, -{Stealth}}
}

\title{Equazioni modulari e sistemi}
\author{Marco Cerbella}
\date{Lezione 16 -- 10 novembre 2026}

\begin{document}
\maketitle

\section{Equazioni modulari: riduzione e insieme delle soluzioni}

Si consideri l'equazione modulare
\[
  ax \equiv b \pmod n, \qquad a, b \in \mathbb{Z},\ n > 0,\ n \in \mathbb{N},
\]
con incognita $x$. Essa ha soluzione se e solo se $\operatorname{MCD}(a, n) \mid b$.

\begin{proposizione}
Siano $a, b \in \mathbb{Z}$, $n > 0$, e $d = \operatorname{MCD}(a, n) \mid b$. Allora l'equazione
\[
  ax \equiv b \pmod n \tag{1}
\]
è equivalente a
\[
  \frac{a}{d}\,x \equiv \frac{b}{d} \pmod{\frac{n}{d}} . \tag{2}
\]
In particolare, se $\bar{x}$ è soluzione di (1) (o di (2)), allora l'insieme delle soluzioni è
\[
  S = \left[\bar{x}\right]_{n/d} = \left[\bar{x}\right]_n \cup \left[\bar{x} + \tfrac{n}{d}\right]_n \cup \left[\bar{x} + 2\tfrac{n}{d}\right]_n \cup \dots \cup \left[\bar{x} + (d-1)\tfrac{n}{d}\right]_n = \left\{\bar{x} + k\tfrac{n}{d} \;\middle|\; k \in \mathbb{Z}\right\}.
\]
\end{proposizione}

In pratica, se $d \neq 1$ si divide l'equazione per $d$ ottenendo un'equazione modulo $n/d$ in cui il coefficiente di $x$ è invertibile (perché $\operatorname{MCD}\!\left(\tfrac{a}{d}, \tfrac{n}{d}\right) = 1$); la si risolve moltiplicando per l'inverso, e poi si descrivono le soluzioni modulo $n$.

\begin{esempio}[$36x \equiv 24 \pmod{105}$]
Si vogliono trovare tutte le soluzioni intere comprese fra $-100$ e $100$.

\emph{Condizione di compatibilità:} $\operatorname{MCD}(36, 105) = 3 \mid 24$. Vero. $[36]_{105}$ non è invertibile.

\emph{Riduzione:} si divide l'equazione per $d = 3$:
\[
  \frac{36}{3}\,x \equiv \frac{24}{3} \pmod{\frac{105}{3}} \quad\Rightarrow\quad 12x \equiv 8 \pmod{35}.
\]
Si cerca l'inverso di $[12]_{35}$ in $\mathbb{Z}_{35}$ con Euclide esteso:
\begin{center}
\begin{tabular}{c|r|r|r}
$q$ & $r$ & $s$ & $t$ \\
\midrule
 & $35$ & $1$ & $0$ \\
 & $12$ & $0$ & $1$ \\
$2$ & $11$ & $1$ & $-2$ \\
$1$ & $1$ & $-1$ & $1 + 2 = 3$ \\
$11$ & $0$ & $1 + 11 = 12$ & $-2 - 33 = -35$ \\
\end{tabular}
\end{center}
Da cui $1 = (-1)\cdot 35 + 3\cdot 12$: l'inverso di $[12]_{35}$ in $\mathbb{Z}_{35}$ è $[3]_{35}$. Moltiplicando per $3$:
\[
  3\cdot 12\,x \equiv 8\cdot 3 \pmod{35} \quad\Rightarrow\quad x \equiv 24 \pmod{35},
\]
quindi $S = [24]_{35} = \{24 + 35k \mid k \in \mathbb{Z}\}$ (soluzioni in $\mathbb{Z}_{35}$).

\emph{Soluzioni fra $-100$ e $100$:}
\[
  24,\ 24 + 35 = 59,\ 24 + 2\cdot 35 = 94,\ (94 + 35 = 129 > 100);
\]
\[
  24 - 35 = -11,\ -11 - 35 = -46,\ -46 - 35 = -81,\ (-81 - 35 = -116 < -100).
\]
Le soluzioni tra $-100$ e $100$ sono $\{-81, -46, -11, 24, 59, 94\}$.

In $\mathbb{Z}_{105}$ le soluzioni sono tre classi: $[24]_{105} \cup [24 + 35]_{105} \cup [24 + 70]_{105} = [24]_{105} \cup [59]_{105} \cup [94]_{105}$.
\end{esempio}

\section{Sistemi modulari}

Un sistema modulare è un insieme di equazioni modulari nella stessa incognita, da risolvere simultaneamente. Prima di tutto si verificano le \emph{condizioni di compatibilità} di ciascuna equazione.

\begin{esempio}
\begin{enumerate}
  \item
  \[
    \begin{cases} 3x \equiv 5 \pmod{12} \\ x \equiv 2 \pmod 4 \end{cases}
  \]
  Condizioni di compatibilità: $\operatorname{MCD}(3, 12) = 3 \nmid 5$ (la prima equazione è impossibile), $\operatorname{MCD}(1, 4) = 1 \mid 2$. Il sistema è impossibile.
  \item
  \[
    \begin{cases} x \equiv 3 \pmod 6 \\ x \equiv 0 \pmod 2 \end{cases}
  \]
  Condizioni di compatibilità: $\operatorname{MCD}(1, 6) = 1 \mid 3$, $\operatorname{MCD}(1, 2) = 1 \mid 0$. Dalla seconda equazione $x = 2k$; sostituendo nella prima: $2k \equiv 3 \pmod 6$, che è \emph{impossibile}: $\operatorname{MCD}(2, 6) = 2 \nmid 3$. Il sistema è impossibile.
  \item
  \[
    \begin{cases} 5x \equiv 3 \pmod 6 \\ x \equiv 1 \pmod 5 \end{cases}
  \]
  Condizioni di compatibilità: $\operatorname{MCD}(5, 6) = 1 \mid 3$, $\operatorname{MCD}(1, 5) = 1 \mid 1$: ok.

  Prima equazione: $5x \equiv 3 \pmod 6$, $\left[\tfrac{1}{5}\right]_6 = [5]_6$; moltiplicando per $5$: $x \equiv 15 \equiv 3 \pmod 6$, cioè $x = 6k + 3$.

  Si sostituisce nella seconda equazione: $6k + 3 \equiv 1 \pmod 5$, cioè $k \equiv 1 - 3 = -2 \equiv 3 \pmod 5$, quindi $k = 5h + 3$.

  Allora $x = 6(5h + 3) + 3 = 30h + 18 + 3 = 30h + 21$, cioè
  \[
    S = [21]_{30} = \{21 + 30h \mid h \in \mathbb{Z}\}.
  \]
\end{enumerate}
\end{esempio}

\end{document}
